% GENERATED FILE. Edit canonical lessons, then run npm run generate:books.
% canonical-source-hash: fnv1a64:35dc89cad676550b
% canonical-lessons: TE-C303-mullu, TE-C303-tiga, TE-C303-veduru, TE-C303-ceraku, TE-C303-puja

\chapter{Thorn, Creeper, Bamboo, Sugarcane, Worship}
\label{ch:te-thorn-creeper-bamboo-sugarcane-worship}
\hlchaptermodality{303}

\begin{chapteropening}
\textbf{By the end of this chapter:} \emph{I can say a thorn, a creeper, bamboo, sugarcane and worship.}

Complete the last lesson of chapter 303: I can say a thorn, a creeper, bamboo, sugarcane and worship.
\end{chapteropening}

\section[\texorpdfstring{\te{ముల్లు} (mullu)}{ముల్లు (mullu)}]{\te{ముల్లు} (mullu) — a thorn}

\label{lesson:TE-C303-mullu}



\begin{quote}
Before the new one: say the Telugu for a wolf, then the Telugu for a crocodile.
\end{quote}

\subsection*{You'll want to know: \te{ముల్లు}}
\textbf{\te{ముల్లు}} — \emph{mullu} — \textquotedblleft{}a thorn\textquotedblright{}.

\begin{grammarlens}[title={What you've built}]
One more word for this chapter.
\end{grammarlens}

\subsection*{Your turn}
\begin{itemize}
\raggedright
  \item \emph{Say it:} \emph{mullu}
  \item \emph{Say it:} \emph{mullu}, once more
  \item \emph{Say it:} say \emph{mullu}
  \item \emph{Recall:} say the Telugu for a wolf, then the Telugu for a crocodile, then say \emph{mullu} again
\end{itemize}

\begin{tcolorbox}[breakable,colback=teal!4,colframe=teal!35!black,title={Before you move on}]
What does \te{ముల్లు} mean? (\textquotedblleft{}A thorn\textquotedblright{}.) Say it once more.
\end{tcolorbox}
\section[\texorpdfstring{\te{తీగ} (tīga)}{తీగ (tīga)}]{\te{తీగ} (tīga) — a creeper, a vine}

\label{lesson:TE-C303-tiga}



\begin{quote}
Before the new one: say the Telugu for a crocodile, then the Telugu for a thorn.
\end{quote}

\subsection*{You'll want to know: \te{తీగ}}
\textbf{\te{తీగ}} — \emph{tīga} — \textquotedblleft{}a creeper, a vine\textquotedblright{}.

\begin{grammarlens}[title={What you've built}]
One more word for this chapter.
\end{grammarlens}

\subsection*{Your turn}
\begin{itemize}
\raggedright
  \item \emph{Say it:} \emph{tīga}
  \item \emph{Say it:} \emph{tīga}, once more
  \item \emph{Say it:} say \emph{tīga}
  \item \emph{Recall:} say the Telugu for a crocodile, then the Telugu for a thorn, then say \emph{tīga} again
\end{itemize}

\begin{tcolorbox}[breakable,colback=teal!4,colframe=teal!35!black,title={Before you move on}]
What does \te{తీగ} mean? (\textquotedblleft{}A creeper, a vine\textquotedblright{}.) Say it once more.
\end{tcolorbox}
\section[\texorpdfstring{\te{వెదురు} (veduru)}{వెదురు (veduru)}]{\te{వెదురు} (veduru) — bamboo}

\label{lesson:TE-C303-veduru}



\begin{quote}
Before the new one: say the Telugu for a thorn, then the Telugu for a creeper.
\end{quote}

\subsection*{You'll want to know: \te{వెదురు}}
\textbf{\te{వెదురు}} — \emph{veduru} — \textquotedblleft{}bamboo\textquotedblright{}.

\begin{grammarlens}[title={What you've built}]
One more word for this chapter.
\end{grammarlens}

\subsection*{Your turn}
\begin{itemize}
\raggedright
  \item \emph{Say it:} \emph{veduru}
  \item \emph{Say it:} \emph{veduru}, once more
  \item \emph{Say it:} say \emph{veduru}
  \item \emph{Recall:} say the Telugu for a thorn, then the Telugu for a creeper, then say \emph{veduru} again
\end{itemize}

\begin{tcolorbox}[breakable,colback=teal!4,colframe=teal!35!black,title={Before you move on}]
What does \te{వెదురు} mean? (\textquotedblleft{}Bamboo\textquotedblright{}.) Say it once more.
\end{tcolorbox}
\section[\texorpdfstring{\te{చెరకు} (ceraku)}{చెరకు (ceraku)}]{\te{చెరకు} (ceraku) — sugarcane}

\label{lesson:TE-C303-ceraku}



\begin{quote}
Before the new one: say the Telugu for a creeper, then the Telugu for bamboo.
\end{quote}

\subsection*{You'll want to know: \te{చెరకు}}
\textbf{\te{చెరకు}} — \emph{ceraku} — \textquotedblleft{}sugarcane\textquotedblright{}.

\begin{grammarlens}[title={What you've built}]
One more word for this chapter.
\end{grammarlens}

\subsection*{Your turn}
\begin{itemize}
\raggedright
  \item \emph{Say it:} \emph{ceraku}
  \item \emph{Say it:} \emph{ceraku}, once more
  \item \emph{Say it:} say \emph{ceraku}
  \item \emph{Recall:} say the Telugu for a creeper, then the Telugu for bamboo, then say \emph{ceraku} again
\end{itemize}

\begin{tcolorbox}[breakable,colback=teal!4,colframe=teal!35!black,title={Before you move on}]
What does \te{చెరకు} mean? (\textquotedblleft{}Sugarcane\textquotedblright{}.) Say it once more.
\end{tcolorbox}
\section[\texorpdfstring{\te{పూజ} (pūja)}{పూజ (pūja)}]{\te{పూజ} (pūja) — worship, a puja}

\label{lesson:TE-C303-puja}



\begin{quote}
Before the new one: say the Telugu for bamboo, then the Telugu for sugarcane.
\end{quote}

\subsection*{You'll want to know: \te{పూజ}}
\textbf{\te{పూజ}} — \emph{pūja} — \textquotedblleft{}worship, a puja\textquotedblright{}.

\begin{grammarlens}[title={What you've built}]
One more word for this chapter.
\end{grammarlens}

\subsection*{Your turn}
\begin{itemize}
\raggedright
  \item \emph{Say it:} \emph{pūja}
  \item \emph{Say it:} \emph{pūja}, once more
  \item \emph{Say it:} say \emph{pūja}
  \item \emph{Recall:} say the Telugu for bamboo, then the Telugu for sugarcane, then say \emph{pūja} again
\end{itemize}

\begin{tcolorbox}[breakable,colback=teal!4,colframe=teal!35!black,title={Before you move on}]
What does \te{పూజ} mean? (\textquotedblleft{}Worship, a puja\textquotedblright{}.) Say it once more.
\end{tcolorbox}
